\documentclass[10pt,a4paper]{amsart}
\usepackage{amsmath,amssymb,amsthm,graphicx}
% Retained preamble for the Stacks Project chapter snapshot in this directory.
% The chapter's own preamble.tex is a copy of the shared Stacks Project preamble; the block
% below is sliced MECHANICALLY, line for line, from it. Nothing was added, removed or
% redefined; the wrapper class, the xr-hyper/hyperref wiring and the externaldocument list
% are omitted because the assembler supplies the document class.
%
% lines taken from preamble.tex: 3, 5, 9, 11, 12, 156, 157, 158, ... (51 lines total)
%
% package coverage restored from the chapter's own preamble (2026-10-13): xy, xy-2cell

\usepackage[all]{xy}
\xyoption{2cell}
\UseAllTwocells

\newenvironment{reference}{\comment}{\endcomment}
\newenvironment{slogan}{\comment}{\endcomment}
\newenvironment{history}{\comment}{\endcomment}
\theoremstyle{plain}
\newtheorem{theorem}[subsection]{Theorem}
\newtheorem{proposition}[subsection]{Proposition}
\newtheorem{lemma}[subsection]{Lemma}
\theoremstyle{definition}
\newtheorem{definition}[subsection]{Definition}
\newtheorem{example}[subsection]{Example}
\newtheorem{exercise}[subsection]{Exercise}
\newtheorem{situation}[subsection]{Situation}
\theoremstyle{remark}
\newtheorem{remark}[subsection]{Remark}
\newtheorem{remarks}[subsection]{Remarks}
\numberwithin{equation}{subsection}
\def\lim{\mathop{\mathrm{lim}}\nolimits}
\def\colim{\mathop{\mathrm{colim}}\nolimits}
\def\Spec{\mathop{\mathrm{Spec}}}
\def\Hom{\mathop{\mathrm{Hom}}\nolimits}
\def\Ext{\mathop{\mathrm{Ext}}\nolimits}
\def\SheafHom{\mathop{\mathcal{H}\!\mathit{om}}\nolimits}
\def\SheafExt{\mathop{\mathcal{E}\!\mathit{xt}}\nolimits}
\def\Sch{\mathit{Sch}}
\def\Mor{\mathop{\mathrm{Mor}}\nolimits}
\def\Ob{\mathop{\mathrm{Ob}}\nolimits}
\def\Sh{\mathop{\mathit{Sh}}\nolimits}
\def\NL{\mathop{N\!L}\nolimits}
\def\CH{\mathop{\mathrm{CH}}\nolimits}
\def\proetale{{pro\text{-}\acute{e}tale}}
\def\etale{{\acute{e}tale}}
\def\QCoh{\mathit{QCoh}}
\def\Ker{\mathop{\mathrm{Ker}}}
\def\Im{\mathop{\mathrm{Im}}}
\def\Coker{\mathop{\mathrm{Coker}}}
\def\Coim{\mathop{\mathrm{Coim}}}
\DeclareMathSymbol{\boxtimes}{\mathbin}{AMSa}{"02}
\def\QCohstack{\mathcal{QC}\!\mathit{oh}}
\def\Cohstack{\mathcal{C}\!\mathit{oh}}
\def\Spacesstack{\mathcal{S}\!\mathit{paces}}
\def\Quotfunctor{\mathrm{Quot}}
\def\Hilbfunctor{\mathrm{Hilb}}
\def\Curvesstack{\mathcal{C}\!\mathit{urves}}
\def\Polarizedstack{\mathcal{P}\!\mathit{olarized}}
\def\Complexesstack{\mathcal{C}\!\mathit{omplexes}}
\def\Pic{\mathop{\mathrm{Pic}}\nolimits}
\def\Picardstack{\mathcal{P}\!\mathit{ic}}
\def\Picardfunctor{\mathrm{Pic}}
\def\Deformationcategory{\mathcal{D}\!\mathit{ef}}
% lane-s11: discard bibliographic \begin{reference} blocks entirely (incl. any \cite inside).
\newsavebox{\limrefbox}
\ifcsname reference\endcsname\expandafter\let\csname reference\endcsname\undefined\fi\ifcsname endreference\endcsname\expandafter\let\csname endreference\endcsname\undefined\fi
\def\citation#1{\relax}
\newenvironment{reference}{\setbox\limrefbox=\vbox\bgroup\let\cite\relax}{\egroup\ignorespaces}
\usepackage{cleveref}
\crefname{originallemma}{original lemma}{original lemmas}
\Crefname{originallemma}{Original lemma}{Original lemmas}
\makeatletter
\newlabel{lemma-descend-finite-presentation}{{01ZM}{source}{}{sourcepointer-1}{}}
\newlabel{lemma-descend-finite-presentation@cref}{{[originallemma][0][]01ZM}{[1][0][]source}{}{}{}}
\makeatother
\crefname{originallemma}{original lemma}{original lemmas}
\Crefname{originallemma}{Original lemma}{Original lemmas}
\makeatletter
\newlabel{lemma-descend-flat-finite-presentation}{{04AI}{source}{}{sourcepointer-2}{}}
\newlabel{lemma-descend-flat-finite-presentation@cref}{{[originallemma][0][]04AI}{[1][0][]source}{}{}{}}
\makeatother
\crefname{originallemma}{original lemma}{original lemmas}
\Crefname{originallemma}{Original lemma}{Original lemmas}
\makeatletter
\newlabel{lemma-directed-inverse-system-has-limit}{{01YX}{source}{}{sourcepointer-3}{}}
\newlabel{lemma-directed-inverse-system-has-limit@cref}{{[originallemma][0][]01YX}{[1][0][]source}{}{}{}}
\makeatother
\crefname{originalreference}{original reference}{original references}
\Crefname{originalreference}{Original reference}{Original references}
\makeatletter
\newlabel{algebra-lemma-flat-finite-presentation-limit-flat}{{02JO}{source}{}{sourcepointer-4}{}}
\newlabel{algebra-lemma-flat-finite-presentation-limit-flat@cref}{{[originalreference][0][]02JO}{[1][0][]source}{}{}{}}
\makeatother
\title{Descent of flat finite-presentation modules along limits}
\author{The Stacks Project Authors}
\date{stacks-project:limits.tex:local-snapshot-20261012; GNU FDL 1.2 or later}
\begin{document}
\maketitle
\noindent\textit{Editorial note.} No omissions in the counted span: the reduction is worked and the algebra core is a precise citation of a proved lemma.\par\medskip
\noindent\textit{Allowed standard prerequisites (editorial).}\begin{itemize}
\item Limit-flat finite presentation. Algebra, Lemma flat-finite-presentation-limit-flat
\end{itemize}\medskip
\ifdefined\hypertarget\hypertarget{sourcepointer-1}{}\fi
\noindent\textit{Source reference.} Original lemma 01ZM: sources/stacks-chow/tags.txt, source lines 2599--2599. Cross-chapter standard prerequisite cited by the author inside the counted span; it is used as a pointer and its own proof is not part of this item.\par\medskip
\ifdefined\hypertarget\hypertarget{sourcepointer-2}{}\fi
\noindent\textit{Source reference.} Original lemma 04AI: sources/stacks-chow/tags.txt, source lines 5617--5617. Cross-chapter standard prerequisite cited by the author inside the counted span; it is used as a pointer and its own proof is not part of this item.\par\medskip
\ifdefined\hypertarget\hypertarget{sourcepointer-3}{}\fi
\noindent\textit{Source reference.} Original lemma 01YX: sources/stacks-chow/tags.txt, source lines 2572--2572. Cross-chapter standard prerequisite cited by the author inside the counted span; it is used as a pointer and its own proof is not part of this item.\par\medskip
\ifdefined\hypertarget\hypertarget{sourcepointer-4}{}\fi
\noindent\textit{Source reference.} Original reference 02JO: sources/stacks-chow/tags.txt, source lines 3343--3343. Cross-chapter standard prerequisite cited by the author inside the counted span; it is used as a pointer and its own proof is not part of this item.\par\medskip
\begin{lemma}
\label{lemma-descend-module-flat-finite-presentation}
With notation and assumptions as in
Lemma \ref{lemma-descend-finite-presentation}.
Let $i \in I$.
Suppose that $\varphi_i : X_i \to Y_i$ is a morphism of schemes
of finite presentation over $S_i$ and that $\mathcal{F}_i$ is a
quasi-coherent $\mathcal{O}_{X_i}$-module of finite presentation.
If the pullback of $\mathcal{F}_i$ to $X_i \times_{S_i} S$ is flat
over $Y_i \times_{S_i} S$, then there exists an index $i' \geq i$
such that the pullback of $\mathcal{F}_i$ to $X_i \times_{S_i} S_{i'}$
is flat over $Y_i \times_{S_i} S_{i'}$.
\end{lemma}

\begin{proof}
(This lemma is the analogue of
Lemma \ref{lemma-descend-flat-finite-presentation}
for modules.)
For $i' \geq i$ denote $X_{i'} = S_{i'} \times_{S_i} X_i$,
$\mathcal{F}_{i'} = (X_{i'} \to X_i)^*\mathcal{F}_i$ and similarly
for $Y_{i'}$. Denote $\varphi_{i'}$ the base change
of $\varphi_i$ to $S_{i'}$. Also set $X = S \times_{S_i} X_i$,
$Y =S \times_{S_i} X_i$, $\mathcal{F} = (X \to X_i)^*\mathcal{F}_i$
and $\varphi$ the base change of $\varphi_i$ to $S$.
Let $Y_i = \bigcup_{j = 1, \ldots, m} V_{j, i}$ be a finite affine open
covering such that each $V_{j, i}$ maps into some affine open of $S_i$.
For each $j = 1, \ldots m$ let
$\varphi_i^{-1}(V_{j, i}) = \bigcup_{k = 1, \ldots, m(j)} U_{k, j, i}$
be a finite affine open covering. For $i' \geq i$ we denote
$V_{j, i'}$ the inverse image of $V_{j, i}$ in $Y_{i'}$ and
$U_{k, j, i'}$ the inverse image of $U_{k, j, i}$ in $X_{i'}$.
Similarly we have $U_{k, j} \subset X$ and $V_j \subset Y$.
Then $U_{k, j} = \lim_{i' \geq i} U_{k, j, i'}$
and $V_j = \lim_{i' \geq i} V_j$
(see Lemma \ref{lemma-directed-inverse-system-has-limit}).
Since $X_{i'} = \bigcup_{k, j} U_{k, j, i'}$ is a finite open covering
it suffices to prove the lemma for each of the morphisms
$U_{k, j, i} \to V_{j, i}$ and the sheaf $\mathcal{F}_i|_{U_{k, j, i}}$.
Hence we see that the lemma reduces to the case that $X_i$ and
$Y_i$ are affine and map into an affine open of $S_i$, i.e., we
may also assume that $S$ is affine.

\medskip\noindent
In the affine case we reduce to the following algebra result.
Suppose that $R = \colim_{i \in I} R_i$. For some $i \in I$
suppose given a map $A_i \to B_i$ of finitely presented $R_i$-algebras.
Let $N_i$ be a finitely presented $B_i$-module.
Then, if $R \otimes_{R_i} N_i$ is flat over $R \otimes_{R_i} A_i$,
then for some $i' \geq i$ the module
$R_{i'} \otimes_{R_i} N_i$ is flat over $R_{i'} \otimes_{R_i} A$.
This is exactly the result proved in
Algebra,
Lemma \ref{algebra-lemma-flat-finite-presentation-limit-flat} part (3).
\end{proof}

\end{document}
